\section{Besoins Fonctionnels}
Les besoins fonctionnels décrivent les actions et les services métier que l'application doit impérativement fournir à ses utilisateurs. Pour SmartFit, l'application s'adresse à deux types d'utilisateurs finaux : les \textbf{Sportifs} et les \textbf{Coachs Professionnels}.

\subsection{Pour les Sportifs}
\begin{itemize}
    \item \textbf{Gestion du compte} : L'utilisateur doit pouvoir créer un compte, se connecter et modifier son profil (âge, poids, taille, objectifs).
    \item \textbf{Programmes d'entraînement} : Le système doit générer un programme d'entraînement personnalisé et adaptatif.
    \item \textbf{Coaching par IA (Vision Posturale)} : L'application doit analyser les mouvements du sportif via la caméra en temps réel et valider ses répétitions.
    \item \textbf{Tracking et Suivi} : L'application doit enregistrer les parcours (GPS), les calories brûlées, et afficher un historique de progression sous forme de graphiques.
    \item \textbf{Nutrition} : Le sportif doit recevoir un plan nutritionnel adapté à ses besoins caloriques.
\end{itemize}

\subsection{Pour les Coachs}
\begin{itemize}
    \item \textbf{Gestion des clients} : Le coach doit pouvoir ajouter, visualiser et suivre les progrès de ses athlètes.
    \item \textbf{Création de programmes} : Le système doit permettre au coach de créer et d'assigner des programmes d'exercices à partir d'une bibliothèque.
    \item \textbf{Communication} : Le coach doit pouvoir communiquer avec ses clients (chat, retours sur les entraînements).
\end{itemize}
